\section{Giới hạn của thực nghiệm}\label{sec:04-evaluation-limitations}
Các kết quả trong chương này cần được diễn giải trong phạm vi dữ liệu, thiết bị và điều kiện đo đã sử dụng.

\paragraph{Dữ liệu mô phỏng.}
Các phép đánh giá phân cụm, xếp hạng và điều phối dùng dữ liệu tổng hợp hoặc bán tổng hợp. Chúng cho phép đối chiếu các phương pháp trong những tình huống được kiểm soát, nhưng chưa đại diện đầy đủ cho hành vi báo tin, nhu cầu và nguồn lực cứu hộ ngoài thực địa. Các chỉ số lợi ích, tổn thất và trễ hạn trong mô phỏng không được xem là hiệu quả cứu hộ đã quan sát thực tế.

\paragraph{Tính đại diện của dữ liệu ảnh.}
Tập ảnh Dataset\_Flood chưa bao quát mọi điều kiện chụp tại hiện trường, như ánh sáng yếu, mưa lớn, ảnh rung hoặc góc nhìn bị che khuất. Kết quả trên tập test phản ánh chất lượng với bộ dữ liệu đã chia; cần ảnh hiện trường độc lập để đánh giá khả năng áp dụng rộng hơn.

\paragraph{Thiết bị và số lượt đo.}
Benchmark nhận diện được thực hiện trên bốn mẫu điện thoại, còn kiểm thử luồng xử lý và giao diện gồm hai run trên Samsung S10e. Số thiết bị và phiên đo còn giới hạn, chưa đủ để kết luận về mọi dòng Android hoặc độ ổn định dài hạn. Thời gian xử lý đã ghi cũng chưa đại diện cho mọi mức tải, nhiệt độ và tình trạng pin của thiết bị.

\paragraph{Điều kiện mạng.}
Các phép đo truyền tải định lượng chủ yếu dùng proxy giới hạn băng thông và độ trễ. Kiểm thử trên điện thoại đã bổ sung tình huống tắt/bật kết nối thật, nhưng chưa bao quát biến động vùng phủ sóng, tắc nghẽn và gián đoạn của mạng viễn thông trong thiên tai. Vì vậy, kết quả xác nhận hành vi trong các điều kiện đã thử, chưa thay thế đánh giá mạng ngoài thực địa.
